\documentclass{beamer}
\usepackage[utf8]{inputenc}

\usetheme{Madrid}
\usecolortheme{default}
\usepackage{amsmath,amssymb,amsfonts,amsthm}
\usepackage{txfonts}
\usepackage{tkz-euclide}
\usepackage{listings}
\usepackage{adjustbox}
\usepackage{array}
\usepackage{tabularx}
\usepackage{gvv}
\usepackage{lmodern}
\usepackage{circuitikz}
\usepackage{tikz}
\usepackage{graphicx}

\setbeamertemplate{page number in head/foot}[totalframenumber]

\usepackage{tcolorbox}
\tcbuselibrary{minted,breakable,xparse,skins}



\definecolor{bg}{gray}{0.95}
\DeclareTCBListing{mintedbox}{O{}m!O{}}{%
  breakable=true,
  listing engine=minted,
  listing only,
  minted language=#2,
  minted style=default,
  minted options={%
    linenos,
    gobble=0,
    breaklines=true,
    breakafter=,,
    fontsize=\small,
    numbersep=8pt,
    #1},
  boxsep=0pt,
  left skip=0pt,
  right skip=0pt,
  left=25pt,
  right=0pt,
  top=3pt,
  bottom=3pt,
  arc=5pt,
  leftrule=0pt,
  rightrule=0pt,
  bottomrule=2pt,
  toprule=2pt,
  colback=bg,
  colframe=orange!70,
  enhanced,
  overlay={%
    \begin{tcbclipinterior}
    \fill[orange!20!white] (frame.south west) rectangle ([xshift=20pt]frame.north west);
    \end{tcbclipinterior}},
  #3,
}
\lstset{
    language=C,
    basicstyle=\ttfamily\small,
    keywordstyle=\color{blue},
    stringstyle=\color{orange},
    commentstyle=\color{green!60!black},
    numbers=left,
    numberstyle=\tiny\color{gray},
    breaklines=true,
    showstringspaces=false,
}
\begin{document}

\title 
{7.4.40}
\date{September 21,2025}


\author 
{Kavin B-EE25BTECH11033}






\frame{\titlepage}
\begin{frame}{Question}
Let a circle be given by $2x\brak{x-a}+y\brak{2y-b}=0$, $\brak{a\neq0, b\neq0}$. Find the condition on $a$ and $b$ if two chords,  each bisected by the $X$ axis, can be drawn to the circle from $\brak{a, \frac{b}{2}}$.\\  
\bigskip
\end{frame}



\begin{frame}{Theoretical Solution}
The equation of the circle is ,
\begin{align}
    2x\brak{x-a}+y\brak{2y-b}=0
\end{align}
\begin{align}
    \implies2x^2 - 2ax + 2y^2 - by = 0
\end{align}
\begin{align}
    \implies x^2 + y^2 - ax - \frac{b}{2}y = 0
\end{align}\\
\bigskip
\end{frame}

\begin{frame}{Theoretical Solution}
The chords are drawn from the point $\vec{P}(a, \frac{b}{2})$ , which lies on the circle since\\
\begin{align}
a^2 + \brak{\frac{b}{2}}^2 - a(a) - \frac{b}{2}\brak{\frac{b}{2}} = a^2 + \frac{b^2}{4} - a^2 - \frac{b^2}{4} = 0
\end{align}\\

\end{frame}


\begin{frame}{Theoretical Solution}
Let a chord start at $\vec{P}(a, \frac{b}{2})$ and end at another point on the circle, $\vec{Q}(x_q, y_q)$. The midpoint of this chord is $\vec{M} = \brak{\frac{a+x_q}{2}, \frac{b/2+y_q}{2}}$.\\ 

Since this chord is bisected by the X-axis,
\begin{align}
\frac{b/2+y_q}{2} = 0 \implies \frac{b}{2} + y_q = 0 \implies y_q = -\frac{b}{2}
\end{align}
\end{frame}

\begin{frame}{Theoretical Solution}
Since $\vec{Q}(x_q, -\frac{b}{2})$ must also lie on the circle,
\begin{align}
    x_q^2 + \brak{-\frac{b}{2}}^2 - a(x_q) - \frac{b}{2}\brak{-\frac{b}{2}} &= 0
\end{align}
\begin{align}
    x_q^2 + \frac{b^2}{4} - ax_q + \frac{b^2}{4} &= 0
\end{align}
\begin{align}
    x_q^2 - ax_q + \frac{b^2}{2} &= 0
\end{align}\\

The problem requires that two such chords can be drawn. This implies that the quadrateic equation for $x_q$ must have two distinct real roots.\\

\end{frame}

\begin{frame}{Theoretical Solution}
Any monic quadratic equation of the form $\lambda^2 - S\lambda + P = 0$ can be interpreted as the \textbf{characteristic equation} of a $2 \times 2$ matrix $\vec{A}$, where $S$ is the trace and $P$ is the determinant of the matrix.
\begin{align}
\mydet{\vec{A} - \lambda \vec{I}} = \lambda^2 - \text{tr}(\vec{A})\lambda + \mydet{\vec{A}} = 0
\end{align}
The roots of this equation are the \textbf{eigenvalues} of the matrix $\vec{A}$.

Here,
\begin{itemize}
    \item  $\lambda = x_q$
    \item  $\text{tr}(\vec{A}) = a$
    \item  $\mydet{\vec{A}} = \frac{b^2}{2}$
\end{itemize}
A fundamental property of a $2 \times 2$ matrix is that it has \textbf{two distinct real eigenvalues} if and only if the discriminant of its characteristic equation is positive. This condition, in terms of its trace and determinant, is:
\begin{align}
(\text{tr}(\vec{A}))^2 - 4\mydet{\vec{A}} > 0
\end{align}\\
\end{frame}


\begin{frame}{Theoretical Solution}
With $\text{tr}(\vec{A}) = a$ and $\mydet{\vec{A}} = \frac{b^2}{2}$, the condition becomes:
\begin{align}
    a^2 - 4\brak{\frac{b^2}{2}} &> 0
\end{align}
\begin{align}
    a^2 - 2b^2 &> 0 
\end{align}
\begin{align}
    a^2 &> 2b^2
\end{align}
This is the required condition on $a$ and $b$. 
\end{frame}
\begin{frame}{Plot}
    \centering
    \includegraphics[width=\columnwidth, height=0.8\textheight, keepaspectratio]{figs/fig.png}
\end{frame}

\end{document}
